\documentclass[journal]{IEEEtran}

% ===================
% Paquetes
% ===================
\usepackage[utf8]{inputenc}
\usepackage[T1]{fontenc}
\usepackage{lmodern}
\usepackage[spanish,es-noshorthands,es-tabla]{babel}
\usepackage{amsmath, amsfonts, amssymb}
\usepackage{graphicx}
\usepackage{cite}
\usepackage{xcolor}
\usepackage{array}
\usepackage{multirow}
\usepackage[hyphens]{url}
\usepackage{tikz}
\usetikzlibrary{babel,arrows.meta,positioning,fit,calc}
\usepackage{pgfplots}
\pgfplotsset{compat=1.18}
\usepgfplotslibrary{groupplots}
% Estilo común de las figuras del experimento (pgfplots). Un color y un
% marcador por escenario, así se distinguen también impresas en gris.
% xmode=log va en cada eje y no aquí: pgfplots lo procesa antes de expandir
% los estilos y dentro de uno da "I do not know the key /tikz/xmode".
\definecolor{esc28}{RGB}{0,114,178}    % producción más rápida (2/8 ms)
\definecolor{esc55}{RGB}{213,94,0}     % tasas equilibradas (5/5 ms)
\definecolor{esc82}{RGB}{0,158,115}    % consumo más rápido (8/2 ms)
\pgfplotsset{
  experimento/.style={
    width=\linewidth, height=4.6cm,
    xtick={1,8,64}, xticklabels={1,8,64},
    xmin=0.7, xmax=90,
    xlabel={Capacidad del buffer (mensajes)},
    tick label style={font=\footnotesize},
    label style={font=\footnotesize},
    legend style={font=\scriptsize, draw=none, fill=none},
    grid=major, grid style={black!12},
    error bars/y dir=both, error bars/y explicit,
    error bars/error mark options={rotate=90, mark size=2pt},
  },
  e28/.style={color=esc28, mark=*, mark size=1.6pt, thick},
  e55/.style={color=esc55, mark=square*, mark size=1.6pt, thick},
  e82/.style={color=esc82, mark=triangle*, mark size=2pt, thick},
}

\usepackage{booktabs}
\usepackage{algorithm}
\usepackage[noend]{algorithmic}
% [H] fija tablas y algoritmos donde se escriben, entre su párrafo de
% introducción y su explicación.
\usepackage{float}
% Sin esto IEEEtran estira los espacios verticales para llenar la columna
% cuando un float [H] pasa a la siguiente, y quedan huecos grandes.
\raggedbottom
% Pseudocódigo en español.
\floatname{algorithm}{Algoritmo}
\renewcommand{\algorithmicif}{\textbf{si}}
\renewcommand{\algorithmicthen}{\textbf{entonces}}
\renewcommand{\algorithmicelse}{\textbf{si no}}
\renewcommand{\algorithmicelsif}{\algorithmicelse\ \algorithmicif}
\renewcommand{\algorithmicendif}{\textbf{fin si}}
\renewcommand{\algorithmicwhile}{\textbf{mientras}}
\renewcommand{\algorithmicdo}{\textbf{hacer}}
\renewcommand{\algorithmicendwhile}{\textbf{fin mientras}}
\renewcommand{\algorithmicrepeat}{\textbf{repetir}}
\renewcommand{\algorithmicuntil}{\textbf{hasta}}
\renewcommand{\algorithmiccomment}[1]{\hfill$\triangleright$~#1}
\usepackage[breaklinks=true,hidelinks]{hyperref}

% Bloque de "pendiente" bien visible en el PDF, para las secciones todavía
% sin escribir. Se reemplaza por contenido real antes de la entrega: el PDF
% entregado no debe contener ninguno (buscar "PENDIENTE" en main.log).
\newcommand{\pendiente}[1]{%
  \begin{quote}\small\color{black!60}%
  \textbf{[PENDIENTE --- #1]}
  \end{quote}%
  \typeout{PENDIENTE: #1}%
}

% Nombres de código en el texto.
\newcommand{\code}[1]{\texttt{#1}}

% IEEEtran pone "Index Terms" en inglés.
\renewcommand{\IEEEkeywordsname}{Palabras clave}

% ===================
% Título y autores
% ===================
% IEEEtran usa \Huge para el título; con este nombre largo ocupaba cuatro
% líneas, así que va con \LARGE (con \huge, "POSIX:" quedaba sola).
\title{\LARGE Buffer Acotado entre Procesos con Memoria Compartida POSIX:\\
Sincronización, Cierre Ordenado y Evaluación Experimental}

\author{
    Isaac Francisco Moreno Fuentes,~\\
    Emmanuel Barrantes Vargas\\
    \thanks{Isaac Francisco Moreno Fuentes y Emmanuel Barrantes Vargas son estudiantes del curso MC 6004 Sistemas Operativos Avanzados, Tecnológico de Costa Rica, II semestre 2026.}
}

% ===================
% Documento
% ===================
\begin{document}
\maketitle

\begin{abstract}
Se implementa un buffer acotado compartido entre procesos con espacios de
direcciones independientes, sobre memoria compartida POSIX. El protocolo usa
un mutex y cuatro variables de condición que viven dentro del segmento,
configurados para compartirse entre procesos, con el predicado de espera de
cada condición declarado y con la escritura de registros fuera de la región
crítica. El cierre es ordenado: un finalizador único corta la admisión
mediante una puerta con contador de referencias, drena lo pendiente, espera
a que todos se retiren y recién entonces destruye las primitivas y libera el
nombre. La verificación combina 205 casos automáticos bajo AddressSanitizer
y UndefinedBehaviorSanitizer con un oráculo que cruza los registros de todos
los procesos y comprueba conservación, unicidad, integridad del contenido y
orden FIFO. Un experimento de 45 corridas ---tres capacidades por tres
relaciones de tasas por cinco repeticiones--- muestra que aumentar la
capacidad del buffer no es una mejora sino un intercambio: eleva el caudal
solo cuando ninguno de los dos lados es el cuello de botella, un 18~\%
entre capacidad~1 y~64 con tasas equilibradas, mientras la latencia sube
siempre, hasta un factor de 66. En régimen saturado la latencia deja de
medir el protocolo y se vuelve la capacidad dividida por la tasa de drenaje.
\end{abstract}

\begin{IEEEkeywords}
buffer acotado, productor-consumidor, memoria compartida POSIX,
variables de condición, cierre ordenado
\end{IEEEkeywords}

% ===================
% Secciones
% ===================
\section{Introducción}
\label{sec:introduccion}

El problema del buffer acotado está resuelto desde que Dijkstra lo planteó
con semáforos~\cite{Dijkstra1965}. Lo que cambia cuando los participantes
son \emph{procesos} con espacios de direcciones independientes, y no hilos
de un mismo programa, es casi todo lo que lo rodea: el estado compartido
tiene que vivir en un objeto con nombre que sobrevive a cada participante,
las primitivas de sincronización pasan a ser datos dentro de ese objeto en
vez de servicios del lenguaje, ninguna dirección virtual vale en otro
proceso, y nadie libera los recursos cuando un proceso termina. Birrell
describe el modelo de hilos del que se parte y menciona esta alternativa
solo para advertir que es más costosa~\cite{Birrell1989}.

Este proyecto implementa ese sistema sobre memoria compartida POSIX. Un
\code{creator} crea una instancia con nombre y capacidad; productores y
consumidores independientes se conectan, se registran, trabajan y se
retiran; un \code{finisher} ejecuta un cierre ordenado que drena lo
pendiente, espera a que todos se retiren y libera los recursos. Un quinto
ejecutable, \code{verify}, cruza los registros de todos los procesos y actúa
como oráculo de corrección. La ejecución ocurre en una sola máquina
GNU/Linux y no cubre la caída arbitraria de un proceso: las garantías se
aplican a participantes que siguen el protocolo y a un cierre normal. Esa
frontera no es una omisión cómoda, y la sección~\ref{sec:pruebas} documenta
qué ocurre exactamente cuando se cruza.

La sección~\ref{sec:arquitectura} describe el segmento compartido y qué dato
es privado de cada proceso; la~\ref{sec:protocolo}, el pseudocódigo del
registro, la inserción, la extracción y el cierre, con sus invariantes y sus
puntos de efecto; la~\ref{sec:pruebas}, las pruebas mínimas y el oráculo;
la~\ref{sec:resultados}, el experimento de 45 corridas;
la~\ref{sec:lecturas}, la relación con Dijkstra y Birrell; y
la~\ref{sec:decisiones}, cinco decisiones donde había una alternativa
razonable que se descartó con evidencia. La sección~\ref{sec:gestion}
documenta la trazabilidad del trabajo.

\section{Arquitectura y representación de memoria}
\label{sec:arquitectura}

\subsection{Ejecutables y módulos}

La compilación produce los cuatro ejecutables del enunciado
(\code{creator}, \code{producer}, \code{consumer} y \code{finisher}) y un
quinto, \code{verify}, que cruza los CSV de una corrida
(Sección~\ref{sec:pruebas}). Cada productor y consumidor es un proceso
independiente que solo conoce el nombre de la instancia. El código común se
divide en los cinco módulos que pide el enunciado y unos pocos de apoyo:

\begin{itemize}
  \item \textbf{Memoria} (\code{memory.c}): crea, valida, mapea y elimina
  el segmento.
  \item \textbf{Buffer} (\code{buffer.c}): arreglo circular puro, sin
  sincronización; sus funciones exigen el mutex ya tomado.
  \item \textbf{Sincronización} (\code{sync.c}): la única capa que toca el
  mutex y las variables de condición; implementa \code{put}, \code{get} y
  la pausa interrumpible.
  \item \textbf{Participantes} (\code{participants.c}): conexión, registro,
  retiro y máquina de estados del cierre.
  \item \textbf{Métricas} (\code{metrics.c}): agregados de la instancia y
  resumen del finalizador.
  \item \textbf{Apoyo}: análisis de argumentos (\code{cli.c}), CSV
  (\code{logger.c}), carga útil verificable (\code{payload.c}), generador
  \emph{xorshift32}~\cite{Marsaglia2003} (\code{rng.c}) y reloj
  monotónico (\code{clock.c}).
\end{itemize}

\subsection{Un segmento por instancia}

Cada instancia es \emph{un único} objeto de memoria compartida POSIX con el
nombre que da el usuario (por ejemplo \code{/soa\_p2\_g1}, visible en Linux
como \code{/dev/shm/soa\_p2\_g1})~\cite{POSIX2024}. No hay semáforos con
nombre ni archivos auxiliares: el mutex y las variables de condición viven
\emph{dentro} del segmento, así que dos instancias con nombres distintos no
comparten ningún recurso. El nombre debe ser \code{/} seguido de 1 a 30
caracteres de \code{[A-Za-z0-9\_-]}.

El segmento contiene un único \code{struct instance}, con el arreglo de
mensajes como miembro flexible al final. La Fig.~\ref{fig:segmento}
muestra su disposición con los desplazamientos medidos en x86-64 con
glibc. La cabecera ocupa 12\,528~bytes y cada mensaje 64, así que el
segmento mide $12\,528 + 64\cdot\mathit{capacidad}$~bytes: 13\,040 con
capacidad 8 y unos 78~KB con 1024.

\begin{figure}[!t]
\centering
\begin{tikzpicture}[
    font=\scriptsize, node distance=0pt,
    campo/.style={draw, text width=5.4cm, align=center, inner xsep=2pt,
                  inner ysep=2.5pt, fill=black!4},
    off/.style={anchor=east, text=black!70},
    prot/.style={anchor=west, text=black!70},
    m/.style={draw, minimum height=0.42cm, inner sep=0pt, font=\tiny,
              fill=black!12, anchor=north west}
  ]
  \node[campo] (c0) {\code{magic} $\cdot$ \code{version} $\cdot$ \code{capacity} $\cdot$ \code{state}};
  \node[campo, below=of c0] (c1) {\code{mutex} (40 B) $+$ \code{not\_full}, \code{not\_empty},
    \code{closing}, \code{all\_left} (48 B c/u)};
  \node[campo, below=of c1] (c2) {\code{head}, \code{tail}, \code{count},
    \code{enqueue\_seq}, \code{dequeue\_seq}, \code{active\_producers},
    \code{active\_consumers}};
  \node[campo, below=of c2] (c3) {\code{refs}: conexiones $+$ bit de puerta cerrada};
  \node[campo, below=of c3] (c4) {\code{producers}: 64 casillas de registro de 64 B};
  \node[campo, below=of c4] (c5) {\code{consumers}: 64 casillas de registro de 64 B};
  \node[campo, below=of c5] (c6) {\code{metrics}: marcas de tiempo, rechazos e
    histograma de latencia (496 cubetas)};
  \node[campo, below=of c6, fill=black!12] (c7) {\code{messages[0..capacidad$-$1]}, 64 B c/u};
  \foreach \i/\o/\p in {0/0/{A, I, I, M}, 1/16/primitivas, 2/248/M, 3/288/A,
                        4/296/M, 5/4\,400/M, 6/8\,504/M, 7/12\,528/M} {
    \node[off]  at (c\i.west) {\o\ };
    \node[prot] at (c\i.east) {\ \p};
  }
  \node[font=\scriptsize\bfseries, anchor=south east] at (c0.north west) {offset};
  \node[font=\scriptsize\bfseries, anchor=south west] at (c0.north east) {protec.};

  % Detalle de un mensaje (no a escala)
  \coordinate (m0) at ($(c7.south west)+(0,-0.45)$);
  \node[m, minimum width=0.72cm] (m1) at (m0) {\code{prod\_id}};
  \node[m, minimum width=0.55cm] (m2) at (m1.north east) {\code{resv}};
  \node[m, minimum width=0.85cm] (m3) at (m2.north east) {\code{sequence}};
  \node[m, minimum width=0.85cm] (m4) at (m3.north east) {\code{enq\_seq}};
  \node[m, minimum width=0.85cm] (m5) at (m4.north east) {\code{enq\_ns}};
  \path let \p1=(m5.east), \p2=(c7.east) in
    node[m, minimum width=\x2-\x1] (m6) at (m5.north east) {\code{payload} (32 B)};
  \draw[densely dotted] (c7.south west) -- (m1.north west);
  \draw[densely dotted] (c7.south east) -- (m6.north east);
  \node[font=\tiny, anchor=north west, align=left, text width=5.6cm, inner xsep=0pt]
    at ($(m1.south west)+(0,-0.05)$)
    {\raggedright payload: \code{sequence} y \code{producer\_id} (copias),
     \code{checksum} \mbox{FNV-1a} y \code{fill[16]}};
\end{tikzpicture}
\caption{Disposición de \code{struct instance} en el segmento (x86-64, glibc
2.39). Protección: A = atómico, I = inmutable tras publicar, M = mutex de
la instancia. Abajo, el detalle de un mensaje de 64~B (no a escala).}
\label{fig:segmento}
\end{figure}

\subsection{Datos compartidos y datos privados}

Todo lo que dos procesos necesitan ver igual está en el segmento: el estado
de la instancia, las primitivas, los índices y contadores del buffer, las
tablas de registro de por vida y los agregados para el resumen final.

Cada proceso guarda en su memoria privada la \emph{dirección} donde mapeó
el segmento (en general distinta entre procesos), el estado de su
generador pseudoaleatorio sembrado una vez con \code{--seed}, su
\code{FILE*} de CSV, sus contadores de espera y adquisición (que copia a su
casilla de registro al retirarse) y el mensaje que construye o procesa.
Así, el procesamiento simulado y la escritura del CSV no necesitan el
mutex.

Ningún campo compartido contiene punteros: el buffer usa los índices
\code{head} y \code{tail} módulo la capacidad sobre \code{messages[]}, y
los mensajes son valores de tamaño fijo que se copian con
\code{buffer\_push} y \code{buffer\_pop}. La carga útil es verificable sin
estado extra: $\mathit{fill}[i] = (\mathit{producer\_id} +
\mathit{sequence} + i) \bmod 256$, y un \emph{checksum} FNV-1a de
32~bits~\cite{FNV} sobre la serialización campo por campo de
\code{sequence}, \code{producer\_id} y \code{fill}. El consumidor la
recalcula al extraer y anota el resultado en su CSV.
Los atómicos del segmento (\code{magic} y \code{refs}) solo son válidos
entre procesos si no usan un candado interno, lo que se exige en tiempo de
compilación con \code{ATOMIC\_INT\_LOCK\_FREE == 2}.

\subsection{Creación exclusiva y publicación}

El \code{creator} sigue un orden fijo:
\begin{enumerate}
  \item valida el nombre y la capacidad;
  \item llama a \code{shm\_open} con \code{O\_CREAT|O\_EXCL};
  \item ajusta el tamaño con \code{ftruncate}, reintentando ante
  \code{EINTR};
  \item mapea el segmento con \code{mmap};
  \item inicializa las primitivas con \code{PTHREAD\_PROCESS\_SHARED} y las
  variables de condición con \code{CLOCK\_MONOTONIC};
  \item escribe \code{version}, \code{capacity} y \code{state = OPEN};
  \item por último, escribe \code{magic} con \emph{release}.
\end{enumerate}

\code{O\_EXCL} hace la creación atómica: entre varios \code{creator}
simultáneos gana uno, y los demás reciben \code{EEXIST} y salen
\emph{sin tocar} el segmento ajeno. Si falla un paso posterior, el
\code{creator} deshace solo lo que creó: destruye las primitivas
inicializadas, desmapea y hace \code{shm\_unlink} del nombre que acaba de
reservar.

Quien se conecta lee \code{magic} con \emph{acquire} y rechaza un
segmento más chico que la cabecera, un \code{magic} distinto de la
constante (\code{creator} sin terminar o segmento ajeno) y una versión,
capacidad o tamaño incompatibles. El par \emph{release}/\emph{acquire} garantiza que quien observa
\code{magic} observa también todo lo escrito antes. Como el kernel entrega
el segmento en ceros y la constante es distinta de cero, una instancia a
medio crear nunca parece lista.

Un mutex no puede proteger esta publicación: quien llega no sabe todavía
si el mutex está inicializado.

\subsection{Propiedad de los recursos}

Cada recurso de una instancia tiene un único responsable de crearlo y
otro de liberarlo, y esa asignación es la que evita fugas y dobles
destrucciones. La Tabla~\ref{tab:propiedad} la resume; el detalle de qué
campo protege cada mecanismo está en \code{include/memory.h}.

\begin{table}[H]
\centering
\caption{Propiedad de los recursos de una instancia}
\label{tab:propiedad}
\footnotesize
\begin{tabular}{@{}p{2.3cm}p{2.5cm}p{2.9cm}@{}}
\toprule
Recurso & Crea & Libera \\
\midrule
Nombre del segmento & \code{creator} (\code{O\_EXCL}) & \code{finisher} (\code{shm\_unlink}) \\
Objeto de memoria & \code{creator} (\code{ftruncate}) & el kernel, al último \code{munmap} \\
Mapeo de un proceso & ese proceso (\code{mmap}) & ese proceso (\code{munmap}) \\
Mutex y condvars & \code{creator}, antes de publicar & \code{finisher}, con $\mathit{refs}=1$ \\
Referencia de conexión & \code{instance\_attach} & \code{instance\_detach} \\
Casilla de registro & \code{participant\_\allowbreak register} & nunca: queda \code{LEFT} \\
CSV de cada proceso & ese proceso & ese proceso \\
\bottomrule
\end{tabular}
\end{table}

Casi todos los recursos los crea el \code{creator} o el propio proceso que
los usa, y los libera el \code{finisher} o ese mismo proceso. Hay dos
excepciones. El objeto de memoria no lo libera nadie del programa: el
\code{finisher} solo elimina el \emph{nombre}, y el kernel libera el
objeto cuando el último proceso lo desmapea. Esa propiedad vuelve seguro
el rechazo de una conexión tardía (Sección~\ref{sec:protocolo}). La
casilla de registro, por su parte, no se libera nunca: queda marcada como
retirada para que su id no pueda volver a usarse. Las primitivas se
destruyen solo cuando el \code{finisher} es el único conectado
($\mathit{refs}=1$).

\section{Protocolo de sincronización y cierre}
\label{sec:protocolo}

\pendiente{Isaac, \#26: Pseudocódigo de registro (puerta \code{refs} con CAS y registro), inserción (\code{sync\_put}), extracción (\code{sync\_get}) y cierre (\code{closing\_begin}, espera en \code{all\_left}, \code{closing\_release}). Tabla de predicados de espera por condvar y quién los hace verdaderos. Invariantes y puntos donde cada operación toma efecto. Ausencia de deadlock. Respuestas a las preguntas orientadoras del enunciado. Fuentes: D3, D9, D10, D18, D21, D26.}

\section{Pruebas y verificación}
\label{sec:pruebas}

\pendiente{Isaac, \#26: Tabla de los 13 casos mínimos (comando, esperado, obtenido, conclusión) a partir de \code{tabla\_pruebas.md}. Verificación automática de mensajes con \code{verify} (D16). Mutaciones detectadas. Sanitizers (\code{make test-san}) y la nota de lo que no demuestran (criterio 3 del \#23, D35). Limitaciones del entorno.}

\section{Resultados del experimento}
\label{sec:resultados}

\pendiente{Emmanuel, \#25: Diseño (9 combinaciones, 5 repeticiones, semillas), entorno (CPU, sistema, compilador) y commit medido. Throughput y latencia media con media $\pm$ desviación, p95 y esperas. Al menos dos gráficas (las tres de \code{graficas/} ya compilan). Si aumentar la capacidad mejora, solo acumula o desplaza la espera. Efecto de los logs. Explicación de la variación. Fuentes: \code{docs/analisis\_experimento.md}, D33.}

% Figuras ya listas (pgfplots, datos en graficas/*.dat): moverlas o
% recortarlas según el texto final.
\begin{figure}[htbp]
\centering
\begin{tikzpicture}
\begin{axis}[experimento, xmode=log, log basis x=2,
    ylabel={Throughput (msgs/s)},
    ymin=380, ymax=840,
    legend style={at={(0.5,1.03)}, anchor=south}, legend columns=3]
  \addplot[e28] table[x=cap, y=e28, y error=e28sd] {graficas/throughput.dat};
  \addlegendentry{2/8 ms}
  \addplot[e55] table[x=cap, y=e55, y error=e55sd] {graficas/throughput.dat};
  \addlegendentry{5/5 ms}
  \addplot[e82] table[x=cap, y=e82, y error=e82sd] {graficas/throughput.dat};
  \addlegendentry{8/2 ms}
  % Techo teórico: la tasa del lado más lento, 4000/m msgs/s.
  \addplot[esc55, dashed, thin, forget plot] coordinates {(0.7,800) (90,800)};
  \addplot[black!50, dashed, thin, forget plot] coordinates {(0.7,500) (90,500)};
\end{axis}
\end{tikzpicture}
\caption{Throughput por capacidad y escenario (media de pausa productor/consumidor),
media $\pm$ desviación estándar entre 5 repeticiones. En línea discontinua, el techo
teórico de cada escenario: 500 msgs/s en 2/8 y 8/2, 800 msgs/s en 5/5.}
\label{fig:throughput}
\end{figure}

\begin{figure}[htbp]
\centering
\begin{tikzpicture}
\begin{axis}[experimento, xmode=log, log basis x=2,
    ymode=log, ylabel={Latencia (ms)},
    ymin=0.03, ymax=400,
    legend style={at={(0.5,1.03)}, anchor=south}, legend columns=3]
  \addplot[e28] table[x=cap, y=e28, y error=e28sd] {graficas/latencia_media.dat};
  \addlegendentry{2/8 ms}
  \addplot[e55] table[x=cap, y=e55, y error=e55sd] {graficas/latencia_media.dat};
  \addlegendentry{5/5 ms}
  \addplot[e82] table[x=cap, y=e82, y error=e82sd] {graficas/latencia_media.dat};
  \addlegendentry{8/2 ms}
  \addplot[e28, dashed, mark options={solid}, forget plot] table[x=cap, y=e28] {graficas/latencia_p95.dat};
  \addplot[e55, dashed, mark options={solid}, forget plot] table[x=cap, y=e55] {graficas/latencia_p95.dat};
  \addplot[e82, dashed, mark options={solid}, forget plot] table[x=cap, y=e82] {graficas/latencia_p95.dat};
\end{axis}
\end{tikzpicture}
\caption{Latencia de inserción a extracción, en escala logarítmica. Línea continua:
media $\pm$ desviación estándar entre repeticiones; discontinua: percentil 95
exacto, calculado desde los CSV de los consumidores.}
\label{fig:latencia}
\end{figure}

\begin{figure}[htbp]
\centering
\begin{tikzpicture}
\begin{groupplot}[experimento, xmode=log, log basis x=2,
    group style={group size=2 by 1, horizontal sep=0.9cm, y descriptions at=edge left},
    width=0.56\linewidth, height=4.4cm,
    ymin=-5, ymax=100, xlabel={Capacidad}]
  \nextgroupplot[title={\footnotesize Productores},
                 ylabel={\% del tiempo de vida},
                 legend to name=leyendaesperas, legend columns=3]
  \addplot[e28] table[x=cap, y=e28, y error=e28sd] {graficas/espera_productor.dat};
  \addlegendentry{2/8 ms}
  \addplot[e55] table[x=cap, y=e55, y error=e55sd] {graficas/espera_productor.dat};
  \addlegendentry{5/5 ms}
  \addplot[e82] table[x=cap, y=e82, y error=e82sd] {graficas/espera_productor.dat};
  \addlegendentry{8/2 ms}
  \nextgroupplot[title={\footnotesize Consumidores}]
  \addplot[e28] table[x=cap, y=e28, y error=e28sd] {graficas/espera_consumidor.dat};
  \addplot[e55] table[x=cap, y=e55, y error=e55sd] {graficas/espera_consumidor.dat};
  \addplot[e82] table[x=cap, y=e82, y error=e82sd] {graficas/espera_consumidor.dat};
\end{groupplot}
\end{tikzpicture}

\ref{leyendaesperas}
\caption{Fracción del tiempo de vida que cada rol pasa bloqueado (productores
esperando espacio, consumidores esperando datos), media $\pm$
desviación estándar entre repeticiones. La espera del consumidor incluye la
espera final hasta el cierre.}
\label{fig:esperas}
\end{figure}


\section{Relación con las lecturas}
\label{sec:lecturas}

\pendiente{Emmanuel, \#27: Dijkstra~\cite{Dijkstra1965} (secciones 3.2, 3.3, 4.1 y 4.3): permisos, capacidad y exclusión; por qué importa el orden de las operaciones; no asumir que la suma instantánea de los dos contadores de semáforo coincide con la capacidad si hay operaciones en curso. Birrell~\cite{Birrell1989}: predicados, separación entre exclusión y coordinación, diferencias entre hilos y procesos, y comparación conceptual con semáforos (sin segunda implementación). Identificar las secciones de Birrell usadas. Clase 5 (diapositiva 30) y clase 7 (diapositivas 12 y 13).}

\section{Discusión de decisiones de diseño}
\label{sec:decisiones}

Las decisiones de este proyecto están registradas una por una ---con su
contexto, las alternativas descartadas y la evidencia cuando la hubo--- en
una bitácora versionada junto al código~\cite{Bitacora2026}. Esta sección
discute tres de ellas, las que tenían una alternativa razonable que se
descartó con evidencia y no por preferencia.

\subsection{Primitivas: mutex y condiciones, no semáforos}

\textbf{Decisión.} Un mutex y cuatro variables de condición POSIX
configuradas para compartirse entre procesos. El enunciado permite las dos
vías y deja la elección al grupo, así que esta es la decisión de la que
cuelgan las dos siguientes.

\textbf{Alternativa descartada.} Semáforos POSIX compartidos entre procesos,
que es la construcción original de Dijkstra para este problema y necesita
menos código: dos contadores, un binario de exclusión y ningún ciclo de
predicado.

\textbf{Evidencia.} Tres razones, ninguna de preferencia.

\emph{Con semáforos hay un orden que interbloquea y hay que recordar no
escribirlo}, el de las dos \code{P} del consumidor
(sección~\ref{sec:lecturas}). Es justo lo que el enunciado prohíbe ---no
esperar mientras se conserva el bloqueo que necesita quien puede satisfacer
esa espera--- y con una variable de condición ese error no se puede
escribir.

\emph{El cierre no necesita permisos falsos.} Despertar a todos los
bloqueados es un \code{broadcast} sobre las cuatro condvars, y lo que cada
uno reevalúa al despertar es el estado real. Por la vía de los semáforos hay
que publicar permisos que no corresponden a ningún mensaje ni a ningún
espacio, y el propio enunciado pide entonces distinguirlos de los reales:
una obligación que esta decisión elimina en vez de cumplir.

\emph{La contabilidad queda en una sola variable.} Un \code{count} bajo un
mutex permite evaluar los dos predicados sobre el mismo estado, lo que con
dos semáforos no vale en todo instante (sección~\ref{sec:lecturas}).

\textbf{Lo que costó.} Las primitivas dejan de ser objetos del sistema y
pasan a vivir dentro del segmento, con \code{PTHREAD\_PROCESS\_SHARED} y un
reloj explícito en las condvars; alguien tiene que crearlas y alguien
destruirlas, que es de donde sale el protocolo de cierre con finalizador
único. Un semáforo con nombre lo habría resuelto el sistema operativo. Ese
es el precio que el proyecto paga por las tres razones de arriba.

\subsection{Tamaño del buffer: no hay un valor por omisión razonable}

\textbf{Decisión.} \code{--capacity} es obligatoria y no tiene valor por
omisión: \code{creator} se niega a crear la instancia si no se la dan.

\textbf{Alternativa descartada.} Darle un valor por omisión cómodo ---8,
digamos--- para poder invocar \code{creator} sin pensar en la capacidad, que
es lo que haría cualquiera que asuma que más espacio nunca estorba.

\textbf{Evidencia.} Ese valor cómodo no existe
(sección~\ref{sec:resultados}): subir la capacidad solo mejora el throughput
cuando ninguno de los dos lados es el cuello de botella (+18\,\% en 5/5~ms,
plano en 2/8~ms), mientras que la latencia sube siempre, hasta un factor de
66. En régimen saturado la latencia ni siquiera es una medición: es la
capacidad dividida por la tasa de drenaje. Cualquier valor por omisión
habría sido bueno para una de las dos métricas y malo para la otra, y el
programa no tiene forma de saber cuál le importa a quien lo corre.

\textbf{Consecuencia práctica.} Elegir la capacidad es elegir un punto en un
intercambio entre caudal y latencia, no una mejora gratuita, así que el
programa obliga a elegirlo de forma explícita. Una simulación de eventos
discretos del mismo escenario, sin mutex ni registro
(\code{scripts/modelo.py}), agrega que el beneficio satura en seco: de
capacidad~128 en adelante el throughput no gana nada, así que el 64 del
barrido ya estaba en el plateau y probar capacidades mayores no habría
aportado información.

\subsection{Granularidad del bloqueo: un solo mutex, y la evidencia de que alcanza}

\textbf{Decisión.} Un mutex por instancia, que protege el anillo, los
contadores, las tablas de participantes y las métricas. \code{buffer.h}
asume que ya está tomado.

\textbf{Alternativa descartada.} Bloqueo más fino ---un candado para el
anillo y otro para las métricas, o secciones críticas separadas por
campo--- con el argumento de reducir contención.

\textbf{Evidencia.} La contención no es el costo, así que el bloqueo fino
optimizaría lo que no duele. Adquirir el mutex cuesta menos de
1~\textmu s por operación en casi todas las celdas (0.24--0.63~\textmu s los
productores, 0.35--1.7~\textmu s los consumidores), entre tres y cuatro
órdenes de magnitud por debajo de las pausas y de las esperas por la
contraparte. Dividir el mutex habría multiplicado los estados posibles y los
órdenes de adquisición ---con el riesgo de interbloqueo que eso trae, el
mismo que Dijkstra señala para el orden de las operaciones--- a cambio de un
microsegundo.

Lo que sí es caro quedó identificado al analizar los resultados, y no es el
candado: despertar a un proceso dormido cuesta unos 60~\textmu s en este
entorno (\code{cond\_signal}, el despacho del planificador y la
readquisición del mutex), dos órdenes de magnitud más que tomar el mutex sin
contención. En el escenario 8/2~ms, donde el consumidor se bloquea en 197 de
cada 200 extracciones, eso es lo que domina la latencia media por mensaje.
La conclusión para el diseño es la inversa de la intuición: conviene
preocuparse por \emph{cuántas veces se duerme un proceso}, no por cuánto
dura la sección crítica.

\section{Gestión del proyecto y trazabilidad}
\label{sec:gestion}

El trabajo se organizó en 32 \emph{Issues} con criterios de aceptación
explícitos, 29 de ellos cerrados, y 31 \emph{Pull Requests} integrados
---todos con al menos una revisión del otro integrante y 11 con dos--- sobre
149 commits. Dos \emph{Issues} de diseño precedieron a la implementación:
uno fijó el formato del segmento y el protocolo de inserción y extracción, y
el otro el registro de participantes, el protocolo de cierre y el formato de
los registros. Ninguno de los dos se implementó antes de estar revisado. Las
decisiones se anotaron en la bitácora a medida que se tomaban, cada entrada
en el \emph{Pull Request} que la implementó~\cite{Bitacora2026}, no al final.

\subsection{Responsabilidades}

El reparto siguió los módulos, no las líneas.

\textbf{Emmanuel Barrantes} trabajó el formato del segmento y el protocolo
de inserción y extracción, el buffer circular, la capa de sincronización, el
reloj monotónico con el generador pseudoaleatorio y la pausa exponencial,
las métricas con el histograma de latencia, el payload verificable, el
verificador \code{verify}, el ejecutable \code{producer} y la primera mitad
del \code{finisher}, el resumen final, el script del experimento y su
análisis, las pruebas de conservación, escala, espera bloqueante y cierre, y
el objetivo de los sanitizers.

\textbf{Isaac Moreno} trabajó el repositorio con su Makefile y su README, el
registro de participantes junto con el protocolo de cierre y el formato de
los registros, el módulo de memoria compartida con el \code{creator}, el
análisis y validación de argumentos, el registrador CSV, el ejecutable
\code{consumer}, la segunda mitad del \code{finisher} con la espera de
retiradas y la liberación, las pruebas de validación, aislamiento y carreras
de ciclo de vida, y la verificación final del código contra el enunciado.

En este informe, las secciones~\ref{sec:arquitectura} a~\ref{sec:pruebas}
son de Isaac y las~\ref{sec:resultados} a~\ref{sec:decisiones} de Emmanuel;
el resto es conjunto.

\subsection{Issues, ramas y revisiones}

Cada cambio entró por una rama con nombre derivado de su \emph{Issue} y un
\emph{Pull Request} que declara la validación ejecutada y queda a revisión
del otro integrante. Los commits pertenecen a quien hizo el trabajo; no se
reescribió historia compartida, no se sustituyó el repositorio al final y
ningún cambio entró como un único commit.

Las revisiones no fueron de forma. Una encontró que la fila
\code{lock\_acquire} del resumen declaraba incluir la readquisición del
mutex al despertar de la espera, cuando el tiempo de bloqueo se mide una
sola vez alrededor de la toma inicial y esa readquisición queda
contabilizada en el tiempo de espera. Los datos lo confirmaron ---la nota
implicaba unos 60~\textmu s y lo medido son 0.89 a 1.74~\textmu s--- y se
corrigió la nota junto con la prueba que la aseveraba. Lo detectó el cruce
de la documentación con el código, no una prueba, que habría pasado igual.

Dos incidentes del flujo quedaron en el historial en vez de borrarse: un
\emph{Pull Request} integrado contra la rama equivocada, que dejó su trabajo
fuera de la rama principal y hubo que rehacerlo contra ella, y una colisión
de numeración en la bitácora cuando dos ramas reclamaron la misma entrada a
la vez.

\section{Conclusiones}
\label{sec:conclusiones}

El protocolo cumple las garantías que el enunciado pide, y lo sostienen dos
clases independientes de evidencia: un verificador externo a la
implementación, que no halló violaciones de conservación, unicidad,
integridad ni FIFO en ninguna corrida, y una suite con intercalados
difíciles de provocar a mano ---cierre en plena actividad, incorporaciones
durante el cierre, cuatro finalizadores compitiendo--- bajo sanitizers sin
reportes. Las dos son necesarias: los sanitizers no buscan carreras y el
verificador no detecta errores de memoria.

Del experimento, el resultado central es que subir la capacidad del buffer
no es una mejora sino un intercambio: aumenta el caudal solo cuando ninguno
de los dos lados es el cuello de botella, mientras la latencia sube siempre
y, en régimen saturado, deja de medir el protocolo para volverse la
capacidad dividida por la tasa de drenaje.

Y queda nítida la frontera entre lo que garantiza el algoritmo ---exclusión,
conservación, orden y terminación--- y lo que concede el sistema operativo,
que decide quién de los que esperan avanza y cuándo. Ninguna propiedad de
corrección depende de esa frontera, pero varios números sí. El alcance
tampoco se amplía al final: no hay recuperación ante la caída arbitraria de
un proceso, no existe herramienta que busque carreras entre procesos
separados sobre memoria compartida, y las mediciones vienen de un binario
sin optimizar sobre un hipervisor.

\section*{Declaración de Transparencia, Uso de IA y Material Reutilizado}

\pendiente{Ambos, como en el P1: uso de Claude Code como herramienta
auxiliar y verificación del contenido contra el código y las corridas
reales. Material reutilizado y su licencia: FNV-1a (dominio público),
\texttt{xorshift32}~\cite{Marsaglia2003}, plantilla \texttt{IEEEtran}.}

\section*{Entrega}

\begin{itemize}
\item \textbf{Repositorio:} \url{https://github.com/ebarrantes07/SOA-P2-GrupoIE}
\item \textbf{Etiqueta:} \texttt{p2-entrega}
\item \textbf{Commit:} \pendiente{hash completo del commit etiquetado.
Bloqueado: la etiqueta \texttt{p2-entrega} todavía no existe. Se llena al
crearla, junto con el README (\#28).}
\end{itemize}

La trazabilidad del trabajo está en la sección «Gestión y trazabilidad» del
README de ese repositorio.


% ===================
% Referencias
% ===================
\bibliographystyle{IEEEtran}
\bibliography{referencias}

\end{document}
